% r03_h1s_a: proof of the two-flow comparison bound (Theorem rot:thm:H1S), part 1: the all-C form, the slices, the elements, degrees per ordering, common moments, the current vertex.

\subsubsection*{Proof of the two-flow comparison bound}

The proof of \Cref{rot:thm:H1S} occupies the rest of this part. Its steps: the lines of every family are written as flow-independent lattice functions sliced by the heat kernel, plus coarse elements (this paragraph); degrees are read per ordering of the scale-summed elements, so that only the connected components of initial segments matter (\Cref{rot:h1s:MS}); each flow's expression is put in a normal form and one subtraction rule is applied, in each flow's own frame, to every component of degree $\le0$ that forms (\Cref{rot:h1s:rule}); after the rule every component containing the weakened element has degree $\ge\min(\alpha,1)$ (\Cref{rot:h1s:Dp}); every element type obeys a two-flow difference law with a positive exponent (\Cref{rot:h1s:P,rot:h1s:E,rot:h1s:U}); and the telescoping over the elements, together with the scale sums, gives the bound with the exponent $\theta_1$ (\Cref{rot:h1s:T} and the proof at the end).

\paragraph{The all-$C$ form and the slices.}
After the constant-background reduction (\Cref{rot:lem:CB}; at $e=0$ nothing is to be reduced) every line of a family is an $A=0$ line of its flow on the infinite lattice: $G^F_r=C^F-\Gamma^F$ (\Cref{rot:lem:Gam}(i)), where $C^F$ is the free propagator of the flow's lattice, $C^X(x,y)=C(x-y)$ and $C^Y(x'',y'')=C(R^{-1}(x''-y''))$ with \emph{one} lattice function $C=(-\Delta_{\eta'}+m_r^2)^{-1}$ on $\eta'\Z^3$ (for the photon, $D=(-\Delta_{\eta'}+\lmu^2\delta_r^2)^{-1}\otimes\mathbf 1_3$). Expanding every line of every graph in $C$ and $\Gamma$ gives \emph{$C/\Gamma$ monomials}, an exact identity; in each monomial the $C$-lines are the same lattice function evaluated in each flow's own frame, the vertices are matched by \Cref{rot:lem:vert} and the counterterm numbers are the same (\Cref{rot:lem:ctsame}), while $\Gamma^F$, $H^F$, $B^F=(A^F)^{-1}$ and the external fields depend on the flow. Write $C$ as the resolvent of the continuous-time simple random walk on $\eta'\Z^3$ with rate $\eta'^{-2}$ per direction, $C(u)=\eta'^{-3}\int_0^\infty e^{-m_r^2s}P_s(u)\,ds$, and slice it in the time variable:
\[
C^{(0)}\defeq\eta'^{-3}\!\int_0^{\ell_1^2},\qquad C^{(j)}\defeq\eta'^{-3}\!\int_{\ell_j^2}^{\ell_{j+1}^2}\ (1\le j\le r-1),\qquad C^{\rm top}\defeq\eta'^{-3}\!\int_1^\infty,
\]
all of $e^{-m_r^2s}P_s(u)\,ds$ ($\ell_r=1$). Gaussian and Poisson tail bounds for the walk give, for every $n\ge0$,
\begin{equation}\label{eq:h1s:B}
|\partial^nC^{(j)}(u)|\le c_n\ell_j^{-1-n}e^{-\delta|u|/\ell_j}\quad(0\le j\le r-1),\qquad|\partial^nC^{\rm top}(u)|\le c_ne^{-c'm_r|u|}(1+|u|)^{-1},
\end{equation}
$\partial$ a lattice difference; a derivative is charged to its line. The decomposition is the same for both flows as lattice functions; the $Y$ flow evaluates it at $R^{-1}(\tau x-\tau y)$, and in $X$-coordinates the $Y$-version obeys \eqref{eq:h1s:B} with $\delta/2$, since $|R^{-1}(\tau x-\tau y)|\ge|x-y|-2c_\tau\eta'$. The heat-kernel slices replace the block-spin slices of \Cref{rot:lem:slice} in the comparison of the all-$C$ lines; both decompositions are exact, and the heat-kernel one is flow-independent as a lattice function, which is what the subtraction rule below needs. (The gain $\eta'^{1/2}$ of \Cref{rot:cor:unif}(a) is recovered, and improved, by \Cref{rot:h1s:P}.)

\paragraph{Elements.}
The \emph{elements} of a monomial, with their per-flow bounds in the unit scale of level $r$, are: (a) the fine slice lines $C^{(j)}$, $D^{(j)}$, $0\le j\le r-1$, with their lattice derivatives, bounded by \eqref{eq:h1s:B}; (a$'$) the top pieces; (b) the coarse elements $\Gamma^F$, $H^F$, $B^F$ (\Cref{rot:lem:Gam}, \Cref{rot:lem:H}), bounded and decaying at the coarse mass, with one derivative and an $\alpha$-H\"older gradient (\cite[Thm.~3.3]{K1} for the standard flow, \Cref{rot:prop:BD}(a) for the hybrid flow); (c) the external fields $E^F\in\{H^F\phi_r,H^F_AA_r$, their covariant derivatives, the Cameron--Martin background$\}$, bounded by the field norms \cite[(3.2)]{K1} on $\supp\chi_r$ and $C^{1,\alpha}$ in the same sense; (d) the vertex functions $w$ (couplings, partitions of unity, the background), smooth with $\|\nabla w\|\le C$ (\Cref{rot:lem:vert}(ii)); (e) at $e>0$ the contour elements of \Cref{rot:prop:BD}(b): contour lines \cite[(3.27)]{K1}, double contour lines, contour vertices \cite[(3.24)--(3.26)]{K1} and the constant-background phases \cite[(3.46)]{K1}; (f) the counterterm numbers $\dm$, $E_1$, the same for both flows. Their matched differences, $X$-version at $x$ against $Y$-version at $\tau x$: (b) $|\Gamma^X-\Gamma^Y\circ\tau|\le C_\Gamma\eta^{1/2}$ (\Cref{rot:lem:Gam}(iii)), $|B^X-B^Y|\le0.2\eta^2$ (\Cref{rot:prop:fc}), coarse slices $\le C\eta^{1/2}\ell_j^{-1}e^{-\delta|x-y|/\ell_j}$ (\Cref{rot:cor:unif}(b)); (c) $|E^X(x)-E^Y(\tau x)|\le C\eta^{1/2}\|E\|$ (\Cref{rot:lem:H}, decaying form, applied to the lift); (d) $|w(x)-w(\tau x)|\le c_\tau\eta'\|\nabla w\|$; (e) contour lines: relative gain $(\eta'/\ell_j)^{1/2}$ on slice $j$, also with one lattice derivative; contour vertices: relative gain $\eta^{1/2}$ against their size $Ce_rp_r^2$; phases: $\le e\,p_r\lmu^{-1}c_\tau\delta_r^{-1/2}\eta'$ (\Cref{rot:prop:BD}(b)); the prefactor $p_r^2\delta_r^{-1/2}$ of the theorem comes from here and from nowhere else. Gradients of (b) and (c) are compared in \Cref{rot:h1s:E}.

\paragraph{The structures.}
With the vertices \cite[(3.18)--(3.23)]{K1} at $B=0$ (\Cref{rot:lem:CB}(1)), the connected all-$C$ line sets of non-positive degree that occur are the following (the list is proved in \Cref{rot:h1s:CL}). SE: the order-$e^2$ scalar self-energy, lines $\{C(x,y),D(x,y)\}$, two current vertices, both derivatives internal. VP: the order-$e^2$ vacuum polarisation, lines $\{C(x,y),C(x,y)\}$, two derivatives on the two lines, together with its seagull tadpole and derivative tadpole (\cite[(3.22)]{K1} with $m=2$), which enter only through the Ward sum of \Cref{rot:h1s:VP}. TAD: the $\lambda$-tadpole $C^{(j)}(0)$ and the seagull scalar tadpole $D^{(j)}(0)$, numbers, identical for both flows and cancelled exactly, slice by slice, by their counterterm shares (\cite[(3.33)]{K1}, \Cref{rot:lem:ctsame}). T: the order-$e^3$ triangle $\{C(a,b),C(b,c),D(a,c)\}$, three current vertices. SC: the order-$e^3$ seagull--current loop $\{C(s,b),D(s,b)\}$. W: the one-particle-irreducible two-point structures of order four: scalar (the $\lambda^2$ sunset; the $e^4$ rainbow, crossed, insertion and seagull graphs) and photon (the $e^4$ self-energy insertion, vertex correction, seagull graphs and the double seagull). L24: Ba{\l}aban's graph \cite[(2.4)]{Ba3}, a single scalar line between two bare current vertices carrying both their derivatives, $\partial\partial C$, with two external photon legs. CT: the counterterm share $c_\sigma$ of a self-energy or tadpole at slice assignment $\sigma$ (a one-site insertion, $|c_\sigma|\le C\ell_{\sigma_{\max}}^{-1}$, a loop line of dimension $-1$ at scale $\sigma_{\max}$), and the order-four share of dimension $0$. The $e^2\lambda$ two-point members are numbers or vanish identically; structures with an odd number of photon legs vanish by charge conjugation; the $(2,1)$-structure at order $e\lambda$ vanishes; four-leg loops have positive degree; contour vertices carry the explicit weight $\ge\frac12$ of \cite[(3.30)]{K1} at their site.

\paragraph{Degrees.}
For a connected set $H$ of elements (lines with their endpoints; a compound, defined in \Cref{rot:h1s:rule}, counts with the vertices and line dimensions it contains) put
\[
D(H)\defeq\sum_{l\in H}d_l+\sum_{v\in V(H)}a_v+3(|V(H)|-1),\qquad d_l=-1-(\text{number of lattice derivatives on }l),
\]
$a_v$ the vertex weight: an explicit factor $\eta'^a$, $a\ge0$, at $v$ counts as $\ell_{j(v)}^a$ with $j(v)\defeq\min\{j_l:l\ni v\}$; a derivative acting on an external field costs nothing, since the externals have bounded gradients. This is Ba{\l}aban's degree \cite[(2.1)--(2.3)]{Ba3} with every derivative charged to its line. The \emph{bare degree} $D_0(H)$ is the same count with the derivative placement of the normal form (NF2) below and before any derivative is transferred. All-$C$ degrees are integers; the only non-integer degrees are vertex weights $\alpha\in(0,1)$ produced by second-order subtractions with an external leg (\Cref{rot:h1s:U}(i)), so every degree is $n+m\alpha$ with $m\ge0$.

\begin{lemma}[the multiscale bound per ordering]\label{rot:h1s:MS}
Let $G$ be a connected set of elements of a monomial, each scale-summed element $l$ carrying the bound $c_l\ell_{j_l}^{d_l}e^{-\delta|x_l-y_l|/\ell_{j_l}}$ at its scale $j_l$ (fixed elements carry a bounded decaying kernel and impose no condition), and fix one ordering class of \cite[(2.7)]{Ba3}: an ordering $l(1),\dots,l(m)$ of the scale-summed elements with the scale sums restricted to $j_{l(1)}\le\dots\le j_{l(m)}$, ties assigned to one class. Suppose that every connected component of every initial segment $G_i\defeq\{l(1),\dots,l(i)\}$ has $D\ge0$, and $D\ge\theta$ if it contains a distinguished element $e_0$ whose constant carries the factor $(\eta'/\ell_{j(e_0)})^\theta$; let $n_0$ be the number of such components with $D=0$. Then the sum over the scales in the class and over all vertices but one is $\le C(1+r)^{n_0}\prod_lc_l$, and with the distinguished element $\le C\eta'^\theta(1+r)^{n_0}\prod_lc_l$, $C=C(G)$. Only the connected components of initial segments are tested.
\end{lemma}
\begin{proof}
This is the proof of \cite[Prop.~2.1]{Ba3} read for one ordering class, which is how Ba{\l}aban runs it: the finest element $l(1)$ is summed first, shrinking its line to a point with the exponent $D(G_1)$; the sum over $j_1\le j_2$ is geometric when $D(G_1)>0$ and costs a factor $(1+r)$ when $D(G_1)=0$; the merged vertex receives the weight $D(G_1)$, and the induction continues with $G_2,G_3,\dots$, at each stage testing only ``the zero-degree subgraphs met along the way'', the connected components of the initial segments. The same per-ordering positivity is King's \cite[p.~665]{K1}. With the distinguished element, every component containing it has exponent $D-\theta>0$ for the geometric sums, and the factor $\eta'^\theta\ell_{j(e_0)}^{-\theta}$ is absorbed by the merged vertex weight $\ell_{j'}^{D-\theta}\le\ell_{j(v)}^{D-\theta}$ at every later stage, so $\eta'^\theta$ survives to the end. A compound enters as a line of its stated dimension at the scale of its leg; its internal scale sums are done in \Cref{rot:h1s:U}(i) uniformly in the positions of its outer vertices, so the first step applies to it verbatim.
\end{proof}

\begin{lemma}[common moments]\label{rot:h1s:Z0}
Let $H(x,t)$ be a kernel built in a flow's own frame from $C$, $D$, their lattice derivatives, lattice Taylor remainders, lattice sums, the counterterm numbers and the bond shifts of \Cref{rot:h1s:cv} (no externals, no coarse elements). Then for $a=0,1$,
\[
\sum_{t\in\Lambda_X}H^Y(\tau x,\tau t)\,(\tau t-\tau x)^{\otimes a}=R^{\otimes a}\sum_{t\in\Lambda_X}H^X(x,t)\,(t-x)^{\otimes a},
\]
the indices compared frame-free (\Cref{rot:lem:vert}(iii)): zero moments are the same numbers and first moments the same vectors in both flows. An $\Oh$-invariant moment is isotropic: a rank-one or rank-three tensor vanishes, a rank-two one is $c\,\delta_{\mu\nu}$. On the tori the kernels are periodisations of the infinite-lattice ones and the moments agree up to $O(e^{-cm_r\cdot50MJ})$ (\Cref{rot:thm:V}(b)).
\end{lemma}
\begin{proof}
$H^Y$ is the conjugate of $H^X$ by the isometry $R$ (one lattice function evaluated in two frames) and $\tau:\Lambda_X\to\Lambda_Y$ is a bijection (\Cref{rot:def:match}), so the left side is $\sum_{v\in R\eta'\Z^3}H(R^{-1}v)v^{\otimes a}=R^{\otimes a}\sum_{u\in\eta'\Z^3}H(u)u^{\otimes a}$. The lattice model is $\Oh$-invariant (King's action on $\eps\Z^3$ with bond fields, in midpoint coordinates for bond kernels), so an $\Oh$-invariant tensor of odd rank vanishes under $-\mathbf 1\in\Oh$ and one of rank two is a multiple of $\delta_{\mu\nu}$.
\end{proof}

\begin{lemma}[the current vertex and the placement of its derivative]\label{rot:h1s:cv}
The current vertex \cite[(3.22)]{K1} ($m=1$) on the bond $b=(x,x+\eta'e_\mu)$ is $e_rA_b\,[\phi(b_+)\cdot q\phi(b_-)]/\eta'$ with $q$ antisymmetric. For two fields $a$, $b$ (a leg and an internal line after Wick contraction),
\[
(D_\mu a)(b)\cdot q\,b(b_-)=-a(b_+)\cdot q\,(D_\mu b)(b)+\eta'^{-1}\big[(a\cdot qb)(b_+)-(a\cdot qb)(b_-)\big],
\]
and summed against $A_b$ over the bonds the last term is $-\sum_x(\partial^*_\mu A_\mu)(x)\,(a\cdot qb)(x)$, the first term of \cite[(2.8)]{Ba3}; if the vertex function $g$ at the vertex is position dependent, the second term of \cite[(2.8)]{Ba3}, $\sum_x(a\cdot qb)(x)\sum_\mu(\partial^{\eta'}_\mu g)(x)A_\mu(x-\eta'e_\mu)$, appears as well. Hence the one lattice derivative of a current vertex may be written on either of its two scalar fields, the other being read at the shifted end of the bond, at the cost of a term in which the derivative acts on the vertex's photon field (and one with an explicit factor $\eta'$ from $\partial^{\eta'}g=\eta'\nabla g+O(\eta'^2)$). The identity is exact in each flow's own frame; a shifted evaluation point is matched by $\tau$ with $c_\tau$ replaced by $c_\tau+1$ in every matched-point bound. A current vertex has exactly two scalar fields and one photon field; the seagull $e_r^2|A|^2|\phi|^2$ and the quartic $\lambda_r|\phi|^4$ carry no derivative.
\end{lemma}
\begin{proof}
$a(b_+)=a(b_-)+\eta'(D_\mu a)(b)$ and $b(b_-)=b(b_+)-\eta'(D_\mu b)(b)$ give $a(b_+)\cdot qb(b_-)/\eta'=(D_\mu a)\cdot qb(b_-)+(a\cdot qb)(b_-)/\eta'=-a(b_+)\cdot q(D_\mu b)+(a\cdot qb)(b_+)/\eta'$; subtract, and sum by parts over the bonds $(x,x+\eta'e_\mu)$. For $a=b=\phi$ the remainder is $\phi\cdot q\phi/\eta'=0$, which is why the one-field form of the vertex has no such term; for two distinct fields it does not vanish, since the photon covariance is $\delta_{\mu\nu}$ times a scalar and not transverse.
\end{proof}
